\section{Privacy, Regulatory Planning, and DPIA}
\label{sec:privacy-gdpr}
\label{sec:dpia}

\subsection{Privacy-by-design commitments}

See Box~\ref{box:system-boundaries} for exclusions and advisory-only rules.
Institutional commitments include:

\begin{itemize}[noitemsep]
  \item local processing on school-controlled systems;
  \item structured zone summaries (not identity records) as operational records;
  \item \textbf{human-in-the-loop} review for elevated items;
  \item purpose limited to safety and emergency coordination.
\end{itemize}

These commitments support DPIA preparation; they do \textbf{not} automatically satisfy GDPR or
EU AI Act obligations without institution-specific legal analysis.

\subsection{GDPR and EU AI Act planning (compact)}

\begin{table}[htbp]
  \centering
  \caption{Regulatory planning topics --- documentation support only (not certification).}
  \label{tab:gdpr-ai-planning}
  \small
  \begin{tabular}{@{}p{2.4cm}p{8.6cm}@{}}
    \toprule
    Topic & Pilot planning consideration \\
    \midrule
    Lawful basis / children &
      DPO-led analysis; guardian transparency; heightened care for minors. \\
    Minimization / purpose &
      Zone summaries only; safety coordination; no profiling or marketing reuse. \\
    Storage / integrity &
      Buffer TTL; event retention in DPIA; segmentation, encryption, RBAC, audit logs. \\
    Accountability / rights &
      Policy versioning; training records; subject requests via school DPO. \\
    EU AI Act (planning) &
      Assess possible high-risk decision-support classification; staff transparency;
      human oversight of advisories; logging for conformity if applicable---\textbf{not claimed here}. \\
    Processors / transfers &
      DPAs with research and IT vendors; subprocessors; breach timelines; transfer safeguards. \\
    \bottomrule
  \end{tabular}
\end{table}

\subsection{Transparency for parents, staff, and students}

Schools should publish a \textbf{plain-language notice} (Appendix~\ref{app:governance-templates})
before Phase~2 and update it when zones change.
Recommended contents:

\begin{enumerate}[noitemsep]
  \item which \textbf{zones} are in scope (spaces, not people by name);
  \item what the pilot \textbf{does not do} (no facial recognition, biometrics, or automated lockdown);
  \item who may access summaries and audit logs;
  \item how to contact the DPO; how to request pause or scope review;
  \item signage where legally required; version date on all public materials;
  \item staff briefing before go-live; age-appropriate communication where required locally.
\end{enumerate}

Consent and lawful basis models vary by jurisdiction---the DPO sets the legal route with counsel.

\subsection{When a DPIA is likely required}

Typically indicated for systematic monitoring in schools, technologies affecting children,
AI-assisted safety decision support, or combined audio-visual processing without identity recognition.
Consult supervisory authority and national education-sector guidance.

\subsection{DPIA preparation checklist}

Completion supports DPO and ethics committees; it does \textbf{not} certify GDPR compliance or
replace a formal institution-signed DPIA.

\begin{longtable}{@{}p{0.5cm}p{4.8cm}p{6.5cm}@{}}
\toprule
\# & Item & Evidence / deliverable \\
\midrule
\endhead
1 & Processing description &
  Zones, modalities, event types, retention, recipients. \\
2 & Necessity \& proportionality &
  Rationale vs.\ manual-only awareness; alternatives. \\
3 & No biometrics attestation &
  Architecture review: no FR/biometric modules. \\
4 & No raw media retention &
  TTL configuration; storage audit; forbidden egress. \\
5 & Lawful basis &
  DPO/legal memo on GDPR basis. \\
6 & Children's data safeguards &
  Guardian communication; access restrictions. \\
7 & Risk identification &
  Link to risk register (Table~\ref{tab:top-risks}; Appendix~\ref{app:risk-register}). \\
8 & Mitigation measures &
  Segmentation, RBAC, audit, training, review gates. \\
9 & Consultation &
  Union, works council, parent council as required. \\
10 & Residual risk acceptance &
  Signed by leadership and DPO. \\
11 & Review schedule &
  Triggers: scope change, incident, new modality. \\
12 & Processor agreements &
  DPAs with research and IT vendors. \\
13 & Breach response &
  Integration with school incident response. \\
14 & Data subject rights &
  Procedures via school DPO. \\
15 & Transparency materials &
  Published notices and signage records. \\
\bottomrule
\end{longtable}

\subsection{Institutional review pack}

Recommended bundle: this document; signed governance charter (Appendix~\ref{app:governance-templates});
protocol mapping matrix; trust-zone diagram; transparency notice; training records; ethics approval
(if human-subjects evaluation); evaluation summary (Section~\ref{sec:evaluation-summary}).

After MVP: update DPIA if scope changes; governance lessons learned; extend/modify/terminate decision.
